%! AP Statistics — Unit 2, Lesson 2.4 — Guided Notes
\documentclass[10pt]{article}
\usepackage{apstats-article}
\usepackage{apstats-boxes}

%=============================================================================
\begin{document}

\pageheader{Unit 2, Lesson 2.4}{Guided Notes}
\namedateperiod

\vspace{0.12in}

% ── OBJECTIVE ────────────────────────────────────────────────────────────────
\begin{objectivebox}
\small By the end of this lesson, I will be able to\dots\\[4pt]
\writeline\\[1pt]
\writeline\\[1pt]
\writeline
\end{objectivebox}

\vspace{0.1in}

% ── VOCABULARY ───────────────────────────────────────────────────────────────
\begin{vocabbox}
\termblanklong{Scatterplot}
\termblanklong{Explanatory variable}
\termblanklong{Response variable}
\termblanklong{Positive association}
\termblanklong{Negative association}
\termblanklong{Form (of a scatterplot)}
\termblanklong{Strength (of a scatterplot)}
\end{vocabbox}

\vspace{0.1in}

% ── HOOK ─────────────────────────────────────────────────────────────────────
\begin{hookbox}
\small\textbf{Opening Question:} Your teacher displays four scatterplots without labels.

\vspace{6pt}
Rank them from \textit{most useful for prediction} to \textit{least useful}, and explain why
you chose that order.\\[8pt]
\writeline\\[1pt]\writeline\\[1pt]\writeline

\vspace{4pt}
\textbf{Key idea from the discussion:} The \blank{5cm} of the pattern in a scatterplot
determines how useful it is for making predictions. Today we learn to describe that pattern precisely.
\end{hookbox}

\vspace{0.1in}

% ── WHAT IS A SCATTERPLOT ────────────────────────────────────────────────────
\begin{notesbox}{Step 1 --- What Is a Scatterplot? (UNC-1.S)}
\small

\noindent A \textbf{scatterplot} displays the relationship between \blank{3.5cm} quantitative
variables measured on \blank{3.5cm} individuals.

\vspace{6pt}
\textbf{Axis rule:}
\begin{itemize}[leftmargin=1.5em, itemsep=2pt]
  \item The \blank{3.5cm} variable (explains or predicts) goes on the \blank{1.5cm}-axis.
  \item The \blank{3.5cm} variable (the outcome) goes on the \blank{1.5cm}-axis.
\end{itemize}

\vspace{6pt}
\textbf{Each point} $(x,\, y)$ represents \blank{5.5cm}.

\vspace{6pt}
\textbf{Class example --- Hours of sleep vs.\ reaction time (ms) for 10 students:}

\vspace{6pt}
\renewcommand{\arraystretch}{1.5}
\begin{tabularx}{\linewidth}{@{} l *{10}{C{0.8cm}} @{}}
\rowcolor{sky}
\textbf{Sleep (hr)} & 5 & 6 & 6 & 7 & 7 & 8 & 8 & 8 & 9 & 9 \\
\textbf{RT (ms)}    & 310 & 290 & 280 & 265 & 250 & 240 & 235 & 220 & 210 & 195 \\
\end{tabularx}

\vspace{10pt}
\begin{center}
\begin{tikzpicture}
  \draw[very thin, gray!35] (0,0) grid[xstep=0.5,ystep=0.5] (8,5);
  \draw[->, thick, navy] (0,0) -- (8.4,0)
        node[right, navy, font=\small] {Sleep (hr)};
  \draw[->, thick, navy] (0,0) -- (0,5.5)
        node[above, navy, font=\small] {RT (ms)};
  \foreach \x/\lbl in {1/5, 2/6, 3/7, 4/8, 5/9, 6/10}
    \node[below, font=\scriptsize] at (\x*1.33,0) {\lbl};
  \foreach \y/\lbl in {0/180, 1/210, 2/240, 3/270, 4/300, 5/330}
    \node[left, font=\scriptsize] at (0,\y) {\lbl};
\end{tikzpicture}
\end{center}

\vspace{4pt}
Plot each point from the table above and label two of them with their coordinates.
\end{notesbox}

\vspace{0.1in}

% ── DOFS FRAMEWORK ───────────────────────────────────────────────────────────
\begin{notesbox}{Step 2 --- Describing the Relationship: DOFS (UNC-1.T)}
\small

\noindent To describe a scatterplot fully, use the \textbf{DOFS} framework:

\vspace{6pt}
\begin{tabularx}{\linewidth}{@{} >{\bfseries\color{navy}}l X @{}}
\textbf{D}irection &
  \blank{2.2cm} association: as $x$ increases, $y$ tends to \blank{2cm}.\\
  & \blank{2.2cm} association: as $x$ increases, $y$ tends to \blank{2cm}.\\
  & \blank{2.2cm} association: knowing $x$ gives no useful information about $y$.\\[6pt]
\textbf{F}orm &
  \blank{2cm}: points follow a straight-line pattern.\\
  & \blank{2cm}: points follow a curved pattern.\\
  & \blank{2cm}: no discernible pattern.\\[6pt]
\textbf{S}trength &
  \blank{2.5cm}: points follow the form very closely.\\
  & \blank{2.5cm}: points follow the form somewhat closely.\\
  & \blank{2.5cm}: points are spread loosely around the form.\\[6pt]
\textbf{O}utliers &
  A point that \blank{4.5cm} from the overall pattern.
\end{tabularx}

\vspace{8pt}
\textbf{Complete DOFS description template (write in context):}\\[6pt]
``There is a \blank{3cm} \blank{3cm} \blank{3cm} association between
[explanatory variable] and [response variable] for these [individuals].
[Outlier sentence if applicable.]''

\vspace{8pt}
\textbf{Write the DOFS description for the sleep vs.\ reaction time scatterplot:}\\[6pt]
\writeline\\[1pt]\writeline\\[1pt]\writeline
\end{notesbox}

\vspace{0.1in}

% ── ADDITIONAL EXAMPLES ───────────────────────────────────────────────────────
\begin{notesbox}{Step 3 --- Practice: Describing Four Scatterplots}
\small

\noindent For each scatterplot your teacher displays, write a complete DOFS description.

\vspace{6pt}
\textbf{Scatterplot A} ($x$ = age of car in years, $y$ = resale price in \$1000s):\\[6pt]
\writeline\\[1pt]\writeline

\vspace{6pt}
\textbf{Scatterplot B} ($x$ = outside temperature in °F, $y$ = hot cocoa sales):\\[6pt]
\writeline\\[1pt]\writeline

\vspace{6pt}
\textbf{Scatterplot C} ($x$ = hours of exercise per week, $y$ = resting heart rate):\\[6pt]
\writeline\\[1pt]\writeline

\vspace{6pt}
\textbf{Scatterplot D} ($x$ = shoe size, $y$ = GPA):\\[6pt]
\writeline\\[1pt]\writeline
\end{notesbox}

\vspace{0.1in}

% ── CONNECTIONS ───────────────────────────────────────────────────────────────
\begin{spiralbox}
\small
\begin{multicols}{2}

\textbf{How this connects:}
\begin{itemize}[itemsep=3pt]
  \item Lesson 2.1: explanatory vs.\ response variable determines which axis each goes on.
  \item Lesson 2.5 (next): correlation $r$ will quantify the \emph{direction} and \emph{strength}
        of a \textit{linear} association with a single number.
  \item Lessons 2.6--2.9: the regression line fits a line to the scatterplot.
\end{itemize}

\columnbreak

\textbf{Reflection:}\\
Why is it important to state the \emph{form} of the association (linear vs.\ curved)?
What might go wrong if we skip it?\\[2pt]
\writeline\\[1pt]\writeline\\[1pt]\writeline

\end{multicols}
\end{spiralbox}

\vspace{0.1in}

\begin{notesbox}{Additional Notes}
\vspace{0pt}
\writeline\\\writeline\\\writeline\\\writeline\\\writeline\\\writeline
\end{notesbox}

\end{document}
